% Part: normal-modal-logic
% Chapter: filtrations
% Section: preliminaries

\documentclass[../../../include/open-logic-section]{subfiles}

\begin{document}

\olfileid[fr]{nml}{fil}{pre}

\olsection{Préliminaires}

Les filtrations permettent d'établir que certains systèmes de logique
modale ont la \emph{propriété du modèle fini}. Si une !!{formula} est vraie
(ou fausse) en un monde d'un modèle de la classe considérée, elle l'est
aussi en un monde d'un modèle \emph{fini} de cette classe. Pour en déduire
la décidabilité par une recherche bornée, on utilise une borne effective
sur la taille du modèle et un test effectif d'appartenance des cadres
finis à la classe. Les filtrations sont
définies relativement à des ensembles de !!{formula}s clos par
sous-formules.

\begin{defn}\ollabel{defn:modallyclosed}
  Un ensemble $\Gamma$ de !!{formula}s est \emph{clos par sous-formules}
  s'il contient toute sous-formule de chaque !!{formula} de~$\Gamma$.
  De plus, $\Gamma$ est \emph{modalement clos} s'il est clos par
  sous-formules et si $!A \in \Gamma$ entraîne $\Box!A,
  \Diamond!A \in \Gamma$.
\end{defn}

Par exemple, pour une !!{formula}~$!A$, l'ensemble de toutes ses
sous-!!{formula}s est clos par sous-!!{formula}s. Dans une filtration
par cet ensemble, $!A$ est vraie en la classe~$[w]$ si et seulement si
elle est vraie en~$w$ dans le modèle initial.

L'ensemble des mondes d'une filtration de~$\mModel{M}$ par~$\Gamma$
est l'ensemble des classes de la relation d'équivalence suivante.

\begin{defn}
Soit $\mModel{M} =\tuple{W, R, V}$ et supposons $\Gamma$ clos par
sous-!!{formula}s. Définissons la relation $\equiv$ sur $W$ entre deux
mondes qui rendent vraies les mêmes !!{formula}s de $\Gamma$ :
\[
u \equiv v \quad \text{si et seulement si } \quad
\forall !A \in \Gamma : \mSat{M}{!A}[u] \Leftrightarrow \mSat{M}{!A}[v].
\]
La classe d'équivalence~$[w]_\equiv$ d'un monde~$w$, ou simplement
$[w]$, est l'ensemble des mondes $\equiv$-équivalents à~$w$ :
\[
[w] = \Setabs{v}{v \equiv w}.
\]
\end{defn}

\begin{prop}
  Étant donnés $\mModel{M}$ et $\Gamma$, la relation $\equiv$ ainsi
  définie est une relation d'équivalence : elle est réflexive,
  symétrique et transitive.
\end{prop}

\begin{proof}
  La relation $\equiv$ est réflexive, car $w$ rend vraies exactement
  les mêmes !!{formula}s de~$\Gamma$ que lui-même. Elle est symétrique :
  si $u$ rend vraies les mêmes !!{formula}s de~$\Gamma$ que $v$,
  alors $v$ fait de même par rapport à~$u$. Elle est transitive : si
  $u$ rend vraies les mêmes !!{formula}s de~$\Gamma$ que~$v$, et $v$
  les mêmes que $w$, alors $u$ rend vraies les mêmes !!{formula}s
  de~$\Gamma$ que~$w$.
\end{proof}

Comme toute relation d'équivalence, $\equiv$ partage~$W$ en classes
formant une \emph{partition} : ce sont des sous-ensembles de~$W$ deux
à deux disjoints, et leur réunion est~$W$. Tout $w \in W$ est un !!{element} d'une de ces
classes, à savoir~$[w]$, puisque $w \equiv w$. Les classes $[w]$
recouvrent donc~$W$. Elles sont deux à deux disjointes : si
$u \in [w]$ et $u \in [v]$, alors $u \equiv w$ et $u \equiv v$ ; par
symétrie et transitivité, $w \equiv v$, d'où $[w] = [v]$.
\end{document}
